% from ega3-1-fr.tex
\oldpage[0\textsubscript{III}]{42}
simple associé à $K_{\bullet\bullet}$
\textup{(\hyperref[0.11.6.2-fr]{11.6.2})}, et est
l'aboutissement commun de six suites spectrales birégulières notées
$\mathop{}^{(t)}E$ (avec $t=a,b,a',b',c$ ou $d$), dont les termes $E^2$ sont
donnés par les formules}
\[
\begin{aligned}
{}^{(a)}E^2_{pq}&=L_pT(H_q(K_{\bullet\bullet}))&
{}^{(b)}E^2_{pq}&=H_p(\mathbf{L}_q^{\mathrm{II}}T(K_{\bullet\bullet}))\\
{}^{(a')}E^2_{pq}&=\mathbf{L}_pT(H_q^{\mathrm I}(K_{\bullet\bullet}))&
{}^{(b')}E^2_{pq}&=H_p(L_qT(K_{\bullet\bullet}))\\
{}^{(c)}E^2_{pq}&=\mathbf{L}_pT(H_q^{\mathrm{II}}(K_{\bullet\bullet}))&
{}^{(d)}E^2_{pq}&=H_p(\mathbf{L}_q^{\mathrm I}T(K_{\bullet\bullet}))
\end{aligned}
\]
\end{proposition}

(On rappelle que nous utilisons la notation $F(A_\bullet)$ pour désigner le
complexe des objets $F(A_i)$ pour tout complexe $A_\bullet=(A_i)$; par
exemple $\mathbf{L}_q^{\mathrm{II}}T(K_{\bullet\bullet})$ désigne le complexe
$(\mathbf{L}_q^{\mathrm{II}}T(K_{i,\bullet}))_{i\geq0}$, où
$\mathbf{L}_q^{\mathrm{II}}T(K_{i,\bullet})$ est l'hyperhomologie d'indice $q$ du
foncteur $T$ par rapport au complexe simple $K_{i,\bullet}$.)

\begin{proof}
Désignons par $L_\bullet$ le complexe simple associé à
$K_{\bullet\bullet}$, de sorte que $L_i=\bigoplus_{r+s=i}K_{rs}$, et
posons
$N_{ij}=\bigoplus_{r+s=i}M_{rsj}$;
il est clair que pour chaque $i$, $N_{i,\bullet}$ est une \emph{résolution
projective} de $L_i$ dans $\boldsymbol{C}$; il résulte donc de
(\hyperref[0.11.6.3-fr]{11.6.3}) et
(\hyperref[0.11.6.4-fr]{11.6.4}) que l'on a un isomorphisme fonctoriel
$
\mathbf{L}_\bullet T(L_\bullet)\xrightarrow{\sim} H_\bullet(T(N_{\bullet\bullet}))$;
comme le complexe simple associé au bicomplexe $T(N_{\bullet\bullet})$ est
aussi associé au tricomplexe $T(M_{\bullet\bullet\bullet})$, cela prouve la
première assertion de l'énoncé.

En outre, $\mathbf{L}_\bullet T(L_\bullet)$ est l'aboutissement des deux suites
spectrales d'hyperhomologie (\hyperref[0.11.6.2-fr]{11.6.2}) de $T$
relativement au complexe simple $L_\bullet$, qui ne sont autres que les suites
$\mathop{}^{(b')}E$ et $\mathop{}^{(a)}E$ respectivement.

Considérons maintenant $M_{\bullet\bullet\bullet}$ comme un bicomplexe
$U_{\bullet\bullet}$ avec
$U_{ij}=\bigoplus_{r+s=j}M_{irs}$;
$H_\bullet(T(M_{\bullet\bullet\bullet}))$ est encore l'aboutissement des deux
suites spectrales du bicomplexe $T(U_{\bullet\bullet})$. Or, pour tout $i$,
$M_{i,\bullet\bullet}$ est un bicomplexe vérifiant les conditions de
(\hyperref[0.11.6.3-fr]{11.6.3}) relativement au complexe simple
$K_{i,\bullet}$; donc
$H_q^{\mathrm{II}}(T(U_{i,\bullet}))=
\mathbf{L}_q^{\mathrm{II}}T(K_{i,\bullet}))$, et la première suite spectrale de
$T(U_{\bullet\bullet})$ n'est autre que $\mathop{}^{(b)}E$. D'autre part, pour
tout $r$, $M_{\bullet,r,\bullet}$ est une résolution de Cartan-Eilenberg du
complexe simple $K_{\bullet,r}$; le calcul fait dans (M, XV, 2) montre que
$H_q^{\mathrm I}(T(M_{\bullet,rs}))=T(H_q^{\mathrm I}(M_{\bullet,rs}))$,
d'où $H_q^{\mathrm I}(T(U_{\bullet,j}))=\bigoplus_{r+s=j}T(H_q^{\mathrm I}(M_{\bullet,rs}))$;
autrement dit, le complexe simple
$H_q^{\mathrm I}(T(U_{\bullet\bullet}))$ n'est autre que le complexe simple
associé au bicomplexe
$T(H_q^{\mathrm I}(M_{\bullet\bullet\bullet})$. Or, pour tout $q$,
$H_q^{\mathrm I}(M_{\bullet\bullet\bullet})$ est une résolution \emph{projective} du
complexe simple $H_q^{\mathrm I}(K_{\bullet\bullet})$ dans la catégorie
$\boldsymbol{K}$; appliquant (\hyperref[0.11.6.3-fr]{11.6.3}), on voit que l'on a
\[
H_p^{\mathrm{II}}
\bigl(T(H_q^{\mathrm I}(M_{\bullet\bullet\bullet}))\bigr)
=\mathbf{L}_pT(H_q^{\mathrm I}(K_{\bullet\bullet})),
\]
et on obtient ainsi la suite $\mathop{}^{(a')}E$. Enfin, les suites
$\mathop{}^{(c)}E$ et $\mathop{}^{(d)}E$ s'obtiennent en intervertissant les
rôles des indices $i$ et $j$ dans la définition du tricomplexe
$M_{\bullet\bullet\bullet}$ et en appliquant à ce nouveau tricomplexe les
raisonnements qui précèdent.
\end{proof}

On dit que $\mathbf{L}_\bullet T(K_{\bullet\bullet})$ est l'\emph{hyperhomologie} de
$T$ relative au \emph{bicomplexe} $K_{\bullet\bullet}$.

\begin{remarks}[11.7.3]
\label{0.11.7.3-fr}
\upshape
\begin{enumerate}
\item[{\rm(i)}] Il résulte de (\hyperref[0.11.6.3-fr]{11.6.3}) que
$\mathbf{L}_\bullet T(K_{\bullet\bullet})$ est un \emph{foncteur homologique} dans la
catégorie des bicomplexes $K_{\bullet\bullet}$ de $\boldsymbol{C}$ limités
inférieurement en chacun de leurs degrés.
